% Revize 2026-10-28: obsah, matematika a návaznost; podrobnosti a časy
% viz planovani/REVIZE_PRIPRAV_1_ROCNIK.md.
% Aktuální porovnání se skutečným obsahem místních PDF Realisticky.cz:
% externi_zdroje/realisticky/01_Zakladni_poznatky/02_Ciselne_obory/010210_Absolutni_hodnota_I/010210_Lekce.pdf
% Rozpis: 1_rocnik/tematicke_plany/Rozpis_vyuky_MAT_1_rocnik.tex, hodina 9;
% SPS.pdf, oddíl 1.2. Návaznost: 05 (znaménka), 07 (osa), 08 (intervaly).
% Učebnice: ucebnice/01_Zakladni_poznatky/Ciselne_obory/Číselné obory.tex,
% oddíl Absolutní hodnota, definice a geometrický význam; bez obecných nerovností.
% Realisticky.cz: externi_zdroje/realisticky/01_Zakladni_poznatky/
% 02_Ciselne_obory/010210_Absolutni_hodnota_I/010210_Lekce.pdf, definice,
% Př. 2–4 a práce s osou. Prostudovány tyto části, nikoli mechanicky převzaty úlohy.
% Všechna konkrétní zadání, obrázek a řešení jsou vlastní tvorba.
% Rovnice jen jako jednoduché hledání bodů v dané vzdálenosti;
% systematické řešení rovnic/nerovnic s absolutní hodnotou patří později.
\documentclass[11pt,a4paper]{article}
% Spolecny styl pro ucitelske pripravy MAT (Work/MAT)
% Verze 1.1, 2026-10-23
\NeedsTeXFormat{LaTeX2e}
\ProvidesPackage{priprava}[2026/10/23 v1.1 Sablona ucitelskych priprav MAT]

% Zakladni typografie a cestina
\RequirePackage[T1]{fontenc}
\RequirePackage[utf8]{inputenc}
\RequirePackage[czech]{babel}
\RequirePackage{lmodern}
\RequirePackage{microtype}
\RequirePackage{amsmath,amssymb}
\RequirePackage{xcolor}
\RequirePackage{enumitem}
\RequirePackage[most]{tcolorbox}
\RequirePackage{tabularx,booktabs,array}
\RequirePackage{fancyhdr}
\RequirePackage{needspace}
\RequirePackage{geometry}
\RequirePackage[hidelinks]{hyperref}

\geometry{a4paper,top=17mm,bottom=18mm,left=23mm,right=23mm}
\definecolor{notegray}{gray}{0.39}
\setlength{\parindent}{0pt}
\setlength{\parskip}{2.5pt plus 1pt minus 0.5pt}
\setlist[itemize]{topsep=2pt,itemsep=1pt,parsep=0pt,leftmargin=1.45em}
\setlist[enumerate]{topsep=2pt,itemsep=2pt,parsep=0pt,leftmargin=1.65em}
\setcounter{secnumdepth}{0}
\renewcommand{\arraystretch}{1.12}

% Zahlavi: vlevo rocnik | globalni tema, vpravo cislo: nazev hodiny
% Zapati zustava zcela prazdne.
\pagestyle{fancy}
\fancyhf{}
\renewcommand{\headrulewidth}{0.35pt}
\renewcommand{\footrulewidth}{0pt}
\setlength{\headheight}{14pt}
\newcommand{\nastavPripravu}[4]{%
  \fancyhead[L]{\small #1\enspace\textbar\enspace #2}%
  \fancyhead[R]{\small #3: #4}%
}

% Minimalisticke cernobile boxy. Kazda definice zvlast.
\tcbset{commonlesson/.style={enhanced,breakable=false,colback=white,colframe=black,
  colbacktitle=white,coltitle=black,boxrule=0.45pt,
  arc=0.8mm,outer arc=0.8mm,left=3mm,right=3mm,top=1.6mm,bottom=1.6mm,
  toptitle=1mm,bottomtitle=0.7mm,fonttitle=\bfseries,
  attach boxed title to top left={xshift=2.5mm,yshift=-1.15mm},
  boxed title style={colback=white,colframe=white,boxrule=0pt,left=1mm,right=1mm}}}
\newtcolorbox{zadaniBox}[1][]{commonlesson,title={Zadání},#1}
\newtcolorbox{prikladBox}[1][]{commonlesson,title={Příklad},#1}
\newtcolorbox{definiceBox}[1][]{commonlesson,title={Definice},#1}

% Spolecny box pro zadani a jeho reseni, bez mezery mezi castmi.
% Pouziti:
%   \begin{prikladReseniBox}
%     Zadani ulohy.
%     \tcblower
%     \begin{align*} ... \end{align*}
%   \end{prikladReseniBox}
% Spodni cast ma jemny sedy podklad a delici linku; lze v ni sazet
% vice radku matematiky. Obe casti maji jen jeden spolecny ramecek.
\tcbset{solutionlesson/.style={
  commonlesson,
  enhanced jigsaw,
  breakable,
  colbacklower=black!3,
  segmentation style={solid,draw=black!25,line width=0.35pt},
  fontlower=\small\color{notegray},
  before lower={\noindent\textbf{Řešení:}\enspace},
  leftlower=3mm,rightlower=3mm
}}
\newtcolorbox{prikladReseniBox}[1][]{solutionlesson,title={Příklad},#1}
\newtcolorbox{zadaniReseniBox}[1][]{solutionlesson,title={Zadání},#1}

% Pomucky pro autora pripravy: sedive poznamky; \reseni pro starsi pripravy
\newcommand{\poznamka}[1]{\par\begingroup\small\color{notegray}\itshape
  \textnormal{\textbf{Poznámka pro učitele.}}\enspace #1\par\endgroup}
\newcommand{\reseni}[1]{\par\begingroup\small\color{notegray}
  \textbf{Očekávané řešení:}\enspace #1\par\endgroup}
\newcommand{\blok}[1]{\needspace{5\baselineskip}\subsection*{#1}}
\newcommand{\vyrok}[1]{\uv{#1}}

\nastavPripravu{1. ročník}{Číselné obory}{09}{Absolutní hodnota}
\begin{document}

\poznamka{Orientačně: osa a motivace 4 min, definice s příklady 8 min,
výpočty a vzdálenost dvou bodů 10 min, samostatná práce s rozborem
11 min, vlastnosti a tolerance 7–10 min, shrnutí 2 min;
základ asi 42–45 min. Rezerva 10–15 min k výběru,
domácí práce 15–20 min. Rovnice zde čteme geometricky, nezavádíme algoritmus
rozštěpování obecného výrazu podle nulových bodů.}

\blok{Poloha má znaménko, vzdálenost nikoli}
Bod $-3$ a bod $3$ mají různé polohy, ale od počátku jsou stejně daleko:
tři jednotky. Vzdálenost nemůže být záporná. Znak pro číslo,
které tuto vzdálenost vyjadřuje, jsou dvě svislé čáry.

\begin{definiceBox}
\textbf{Absolutní hodnota} reálného čísla $a$, značená $|a|$,
je vzdálenost jeho obrazu od nuly na číselné ose. Platí
\[
|a|=\begin{cases}a,&a\ge0,\\-a,&a<0.\end{cases}
\]
\end{definiceBox}

Proto $|0|=0$, $|5|=5$ a $|-5|=-(-5)=5$.
Absolutní hodnota je \emph{nezáporná}, nikoli vždy kladná.
Čísla opačná mají stejnou absolutní hodnotu: $|-a|=|a|$.
Samotný zápis $-a$ však nemusí být záporný, jak víme z hodiny 05.

\begin{prikladReseniBox}[title={P1: záleží na místě svislých čar},breakable=false]
Vypočti $|-7|$, $-|7|$, $|-7+2|$, $|-7|+2$ a $|-\frac35|$.
\tcblower
$|-7|=7$, ale $-|7|=-7$.
Dále $|-7+2|=|-5|=5$, zatímco $|-7|+2=7+2=9$.
Nakonec $|-\frac35|=\frac35$.
Absolutní hodnota působí jen na výraz uvnitř svých čar;
například výraz $-|7|$ může být záporný.
\end{prikladReseniBox}

\begin{prikladReseniBox}[title={P2: běžné pořadí operací stále platí},breakable=false]
Vypočti $3-2\cdot|-4+1|+|-8|:2$ a $\bigl||-5|-8\bigr|$.
\tcblower
Nejprve spočítáme výrazy uvnitř čar:
\begin{align*}
3-2\cdot|-4+1|+|-8|:2&=3-2\cdot3+8:2=3-6+4=1,\\
\bigl||-5|-8\bigr|&=|5-8|=|-3|=3.
\end{align*}
U vnořeného zápisu postupujeme od nejvnitřnějších čar ven.
\end{prikladReseniBox}

\poznamka{Vnořené čáry piš na tabuli v různých výškách. Ptej se:
\uv{Který celý výraz se právě nachází uvnitř?} Chybu neopravuj pouhým
smazáním záporných znamének, ale vrať se k hodnotě a vzdálenosti.}

\newpage
\blok{Vzdálenost mezi dvěma body}
\begin{definiceBox}
\textbf{Vzdálenost bodů} se souřadnicemi $a,b\in\mathbb R$ na číselné ose
je $d(a,b)=|a-b|$. Stejně dostaneme $d(a,b)=|b-a|$.
\end{definiceBox}

Je-li $a\ge b$, je vzdálenost $a-b$; je-li $a<b$, je to $b-a$.
Obě situace shrnuje absolutní hodnota rozdílu.

\blok{Základní vlastnosti a jejich význam}
Pro reálná $a,b$ platí
\begin{align*}
|a|&\ge0, & |a|=0&\ \Leftrightarrow\ a=0, &|-a|&=|a|,\\
|ab|&=|a|\,|b|, & \left|\frac ab\right|&=\frac{|a|}{|b|}\quad(b\ne0).
\end{align*}
U součinu a podílu nejprve zjistíme jeho znaménko: absolutní hodnota
odstraní znaménko výsledku a ponechá jeho velikost, stejně jako součin
či podíl velikostí obou čísel. Například
$|(-3)\cdot2|=6=|-3|\cdot|2|$.

Pro součet platí \emph{trojúhelníková nerovnost}
$|a+b|\le|a|+|b|$.
Na ose ji čteme jako dva posuny: konečná vzdálenost od počátku
nepřesáhne součet délek obou posunů. Při stejném směru se délky
sčítají, při opačných směrech se část cesty vracíme.
Například $|5+(-3)|=2<8=|5|+|-3|$; rovnost tedy nemusí platit.

\begin{prikladReseniBox}[title={P3: rozdíl a délka úseku},breakable=false]
Urči vzdálenost bodů $A=-3$ a $B=2$ i bodů $C=-4$ a $D=-1$.
\tcblower
$d(A,B)=|-3-2|=|-5|=5$ a $d(C,D)=|-4-(-1)|=|-3|=3$.
V prvním případě jdeme tři jednotky k nule a další dvě k $B$;
ve druhém jsou oba body vlevo od nuly a dělí je pouze tři jednotky.
\begin{center}
\begin{tikzpicture}[x=1.35cm,y=0.65cm]
\draw[->] (-4.5,0)--(3,0);
\foreach \x in {-4,-3,-2,-1,0,1,2}{\draw (\x,-0.08)--(\x,0.08) node[below=4pt] {$\x$};}
\fill (-3,0) circle (2pt) node[above=3pt] {$A$};
\fill (2,0) circle (2pt) node[above=3pt] {$B$};
\draw[<->] (-3,0.9)--(2,0.9) node[midway,above] {$5$};
\end{tikzpicture}
\end{center}
\end{prikladReseniBox}

\begin{prikladReseniBox}[title={P4: hledej všechny body v dané vzdálenosti},breakable=false]
Najdi všechna reálná $x$, pro která platí a) $|x|=4$;
b) $|x|=0$; c) $|x|=-2$. Zdůvodni, proč další možnosti nejsou.
\tcblower
a) Ve vzdálenosti $4$ od nuly jsou právě dva body: $x=-4$ a $x=4$.
Na každé straně počátku je právě jeden.
b) Nulovou vzdálenost od nuly má pouze $x=0$.
c) Žádné reálné číslo nevyhovuje; vzdálenost není záporná.
Dosazení $|-4|=|4|=4$ a $|0|=0$ potvrzuje nalezené možnosti.
\end{prikladReseniBox}

\begin{prikladReseniBox}[title={P5: tolerance rozměru},breakable=false]
Jmenovitá délka je $10$ mm a přípustná odchylka nejvýše $0{,}2$ mm.
Vyjádři podmínku pomocí absolutní hodnoty a intervalem.
Vyhovují délky $9{,}8$ mm a $10{,}3$ mm?
\tcblower
Vzdálenost skutečné hodnoty $d$ od $10$ musí být nejvýše $0{,}2$:
$|d-10|\le0{,}2$, kde $d$ vyjadřujeme v milimetrech.
Odtud $9{,}8\le d\le10{,}2$, tedy $d\in\langle9{,}8;10{,}2\rangle$.
První délka vyhovuje, protože $|9{,}8-10|=0{,}2$;
druhá nevyhovuje, protože $|10{,}3-10|=0{,}3>0{,}2$.
\end{prikladReseniBox}

\blok{Samostatná práce}
\begin{zadaniReseniBox}[title={S1: výpočet a vzdálenost},breakable=false]
\begin{enumerate}[label=\alph*)]
\item Vypočti $|-12|$, $|0|$, $-|-6|$, $|3-8|$, $|-3|-8$, $|-0{,}4|$.
\item Vypočti $2\cdot|-3|-|4-9|$ a $\bigl|-2+|-7|\bigr|$.
\item Urči vzdálenosti dvojic $(-5,1)$, $(-6,-2)$, $(\frac12,-\frac34)$.
\item Najdi všechna reálná $x$ s $|x|=\frac32$ a s $|x|=-\frac12$.
\end{enumerate}
\tcblower
a) $12$, $0$, $-6$, $5$, $-5$, $0{,}4$.
b) $6-5=1$; $|-2+7|=5$.
c) $6$, $4$, $|\frac12+\frac34|=\frac54$.
d) $x=-\frac32$ nebo $x=\frac32$; druhé podmínce nevyhovuje žádné číslo.
\end{zadaniReseniBox}

\begin{zadaniReseniBox}[title={S2: obhaj nebo vyvrať},breakable=false]
Platí pro všechna reálná $a,b$: a) $|a|>0$;
b) $|a+b|=|a|+|b|$; c) $|a-b|=|b-a|$?
\tcblower
a) Ne, pro $a=0$ je $|a|=0$; správně $|a|\ge0$.
b) Ne, pro $a=2,b=-2$ vlevo $0$, vpravo $4$.
c) Ano, oba zápisy vyjadřují vzdálenost týchž dvou bodů.
\end{zadaniReseniBox}

\poznamka{Na závěr nech slovně rozlišit polohu $-5$ od vzdálenosti $5$.
Pro slabší žáky nahraď vnořený výraz konkrétním měřením na ose.
V příští hodině následuje krátký test a rozbor; připomeň i logiku z hodin 02–04.}

\newpage
\blok{Rezerva: uvažování nad osou}
\begin{zadaniReseniBox}[title={R1: volitelně, 10–15 minut},breakable=false]
\begin{enumerate}[label=\alph*)]
\item Najdi všechny body ve vzdálenosti $3$ od bodu $2$.
Zapiš kontrolu pomocí absolutní hodnoty rozdílu.
\item Která reálná čísla jsou od nuly vzdálena méně než $2$?
Zapiš je nerovností i intervalem. Co se změní pro vzdálenost nejvýše $2$?
\item Najdi bod stejně vzdálený od $-4$ a $2$. Vysvětli jedinečnost.
\item Je-li $a<0$, jaké jsou hodnoty $|a|$, $|-a|$ a $-|a|$?
Ověř odpověď pro $a=-6$.
\end{enumerate}
\tcblower
a) Tři jednotky vlevo a vpravo od $2$: $x=-1$ nebo $x=5$.
Kontroly $|-1-2|=3$ a $|5-2|=3$; na každé straně je právě jeden bod.
b) $-2<x<2$, interval $(-2;2)$; pro nejvýše $2$ je to
$-2\le x\le2$, interval $\langle-2;2\rangle$.
c) Střed úseku je $-1$, od obou krajů má vzdálenost $3$.
Mimo úsek je jeden kraj vždy blíže; uvnitř se obě vzdálenosti rovnají
jen při rozdělení celkové délky $6$ napůl.
d) $|a|=-a$, $|-a|=-a$, $-|a|=a$. Pro $a=-6$ postupně $6,6,-6$.
\end{zadaniReseniBox}

\blok{Domácí procvičování}
\begin{zadaniReseniBox}[title={D1: přibližně 15–20 minut},breakable=false]
\begin{enumerate}[label=\arabic*)]
\item Vypočti $|-15|$, $|2-9|$, $-|2-9|$, $|-\frac78|$.
\item Vypočti $|-6+2|+3\cdot|-2|$ a $\bigl||-9|-4\bigr|$.
\item Vypočti $5-2\cdot|-4|$; vysvětli, proč výsledek může být záporný.
\item Nakresli body $-2$ a $3{,}5$ a urči jejich vzdálenost.
Totéž proveď pro $-\frac32$ a $-\frac14$.
\item Najdi všechna reálná $x$ s $|x|=2{,}5$, s $|x|=0$ a s $|x|=-3$.
\item Vyvrať vhodnými čísly tvrzení $|a-b|=|a|-|b|$ pro všechna reálná
$a,b$. Napiš, co skutečně vyjadřuje levá strana.
\end{enumerate}
\tcblower
1. $15$, $7$, $-7$, $\frac78$.
2. $4+6=10$; $|9-4|=5$.
3. $5-8=-3$; absolutní hodnota zaručuje nezápornost pouze hodnoty uvnitř
příslušných čar, ne celého výrazu s odčítáním.
4. $|-2-3{,}5|=5{,}5$; $|-\frac32+\frac14|=\frac54$.
5. $x=-2{,}5$ nebo $2{,}5$; pouze $x=0$; žádné reálné číslo.
6. Například $a=1,b=3$: vlevo $2$, vpravo $-2$. Levá strana je vzdálenost
bodů se souřadnicemi $a,b$ a nemůže být záporná.
\end{zadaniReseniBox}

\end{document}
